\documentclass[spanish,letterpaper,oneside]{article}
\usepackage[utf8]{inputenc}
\usepackage[spanish]{babel}
\usepackage{geometry}
\geometry{letterpaper, margin=2.5cm}
\usepackage{xcolor}
\usepackage{enumitem}
\usepackage{tcolorbox}

\begin{document}
\begin{titlepage}
    \centering
    \vspace*{2cm}
    {\Huge \textbf{Universidad Autónoma de Coahuila} \par}
    \vspace{1cm}
    {\Large Facultad de Medicina \par}
    \vspace{2cm}
    {\LARGE \textbf{Examen de Ingreso: Bioquímica (Premedicina)} \par}
    \vspace{1cm}
    {\Large Guía Autocontenida y Banco de Reactivos EXANI-II \par}
    \vspace{3cm}
    {\large Elaborado por: Aspirante a Médico Cirujano \par}
    \vfill
    {\large \today \par}
\end{titlepage}

\tableofcontents
\newpage

\section{Módulo Teórico: Bioquímica de Alto Rendimiento}
\subsection{Biomoléculas y Macromoléculas}
Las 4 macromoléculas de la vida:
\begin{itemize}
    \item \textbf{Carbohidratos:} Monosacáridos (glucosa, galactosa, fructosa), disacáridos (sacarosa, lactosa) y polisacáridos (glucógeno en animales, almidón en plantas). Unidos por enlace \textbf{O-glucosídico}.
    \item \textbf{Lípidos:} Insolubles en agua. Triglicéridos (energía y aislamiento), Fosfolípidos (membranas celulares, son anfipáticos), Esteroides (derivados del colesterol). Enlaces \textbf{éster}.
    \item \textbf{Proteínas:} Cadenas de aminoácidos unidos por enlace \textbf{peptídico}. Estructuras: Primaria (secuencia lineal), Secundaria (alfa-hélice, lámina beta por puentes de H), Terciaria (plegamiento tridimensional activo), Cuaternaria (más de una cadena, ej. Hemoglobina).
    \item \textbf{Ácidos Nucleicos:} Polímeros de nucleótidos (fosfato + azúcar + base nitrogenada). ADN (Desoxirribosa; A, T, C, G) y ARN (Ribosa; A, U, C, G). Enlace \textbf{fosfodiéster}.
\end{itemize}

\subsection{Metabolismo Celular (Respiración Aerobia)}
\begin{tcolorbox}[colback=green!5!white,colframe=green!50!black,title=Las 3 etapas de la Respiración Celular Aerobia]
1. \textbf{Glucólisis:} En el citoplasma. Anaeróbica. Convierte 1 Glucosa (6C) en 2 Piruvatos (3C). Ganancia neta: 2 ATP y 2 NADH.\\
2. \textbf{Ciclo de Krebs (Ácido Cítrico):} En la matriz mitocondrial. El piruvato pasa a Acetil-CoA y entra al ciclo. Genera por vuelta: 3 NADH, 1 FADH2, 1 ATP (o GTP), 2 CO2.\\
3. \textbf{Cadena de Transporte de Electrones y Fosforilación Oxidativa:} En la membrana mitocondrial interna. Los electrones de NADH/FADH2 bombean protones creando un gradiente. La \textit{ATP sintasa} produce ATP masivamente (aprox. 34 ATP). El $ es el aceptor final (forma agua). Total: 36-38 ATP por glucosa.
\end{tcolorbox}

\subsection{Enzimología}
\begin{itemize}
    \item \textbf{Enzimas:} Catalizadores biológicos proteicos (disminuyen la energía de activación de una reacción). Se unen a un sustrato específico en su \textit{sitio activo}.
    \item \textbf{Inhibición Competitiva:} El inhibidor se asemeja al sustrato y compite por el sitio activo. Se puede vencer aumentando el sustrato.
    \item \textbf{Inhibición No Competitiva:} El inhibidor se une a un sitio alostérico, cambiando la forma de la enzima e inactivándola.
\end{itemize}

\subsection{Dogma Central de la Biología Molecular}
El flujo de información genética es unidireccional (con excepciones como virus retrotranscriptores):
\begin{enumerate}
    \item \textbf{Replicación (ADN $\rightarrow$ ADN):} Semisconservativa, en el núcleo (fase S). La enzima principal es la \textit{ADN Polimerasa}.
    \item \textbf{Transcripción (ADN $\rightarrow$ ARNm):} En el núcleo. La enzima es la \textit{ARN Polimerasa}. Produce un ARNm inmaduro.
    \item \textbf{Traducción (ARNm $\rightarrow$ Proteína):} En los ribosomas (citoplasma/RER). El ARNt trae el aminoácido leyendo los codones (tripletes).
\end{enumerate}

\newpage
\section{Banco de Reactivos: Simulación EXANI-II}
\begin{enumerate}[label=\textbf{\arabic*.}]
    \item \textbf{Biomolécula compuesta por cadenas de aminoácidos, la cual cumple funciones estructurales, catalíticas (enzimas) e inmunológicas (anticuerpos):}
    \begin{enumerate}
        \item[A)] Carbohidratos
        \item[B)] Lípidos
        \item[C)] Proteínas
        \item[D)] Ácidos nucleicos
    \end{enumerate}
    \textit{Respuesta: C. Las proteínas son los únicos polímeros de aminoácidos y son extremadamente versátiles (ej. enzimas, hemoglobina, queratina).}

    \item \textbf{¿Qué nombre recibe el enlace químico que une a los monosacáridos para formar oligosacáridos y polisacáridos?}
    \begin{enumerate}
        \item[A)] Peptídico
        \item[B)] Glucosídico
        \item[C)] Fosfodiéster
        \item[D)] Éster
    \end{enumerate}
    \textit{Respuesta: B. El enlace glucosídico (u O-glucosídico) une azúcares. El peptídico une aminoácidos, fosfodiéster nucleótidos y éster lípidos.}

    \item \textbf{La sacarosa (azúcar de mesa) es un disacárido formado por la unión de:}
    \begin{enumerate}
        \item[A)] Glucosa + Glucosa
        \item[B)] Glucosa + Galactosa
        \item[C)] Fructosa + Galactosa
        \item[D)] Glucosa + Fructosa
    \end{enumerate}
    \textit{Respuesta: D. Glucosa + Fructosa = Sacarosa. (Glucosa+Glucosa=Maltosa, Glucosa+Galactosa=Lactosa).}

    \item \textbf{Molécula lipídica con una cabeza hidrofílica y dos colas hidrofóbicas, que conforma la base de la bicapa de la membrana celular:}
    \begin{enumerate}
        \item[A)] Triglicérido
        \item[B)] Fosfolípido
        \item[C)] Colesterol
        \item[D)] Esfingolípido
    \end{enumerate}
    \textit{Respuesta: B. Los fosfolípidos son la estructura anatómica primaria de toda membrana celular gracias a su carácter anfipático.}

    \item \textbf{La glucólisis es una ruta metabólica universal que ocurre en el citosol celular. ¿Cuál es el producto final neto (por cada glucosa) al concluir esta vía?}
    \begin{enumerate}
        \item[A)] 2 Piruvatos, 2 ATP, 2 NADH
        \item[B)] 2 Acetil-CoA, 4 ATP, 2 NADH
        \item[C)] 4 Lactatos, 2 ATP
        \item[D)] 2 Piruvatos, 36 ATP, 2 FADH2
    \end{enumerate}
    \textit{Respuesta: A. Se producen 4 ATP en total, pero se gastan 2 en la fase preparatoria, por lo que la ganancia neta es de 2 ATP.}

    \item \textbf{¿En qué organelo celular eucariota se lleva a cabo el Ciclo de Krebs (Ciclo del Ácido Cítrico)?}
    \begin{enumerate}
        \item[A)] Citoplasma
        \item[B)] Ribosomas
        \item[C)] Matriz mitocondrial
        \item[D)] Aparato de Golgi
    \end{enumerate}
    \textit{Respuesta: C. A diferencia de la glucólisis (citosol), el Ciclo de Krebs ocurre dentro del organelo productor de energía, la mitocondria (específicamente en la matriz líquida).}

    \item \textbf{Durante la Cadena de Transporte de Electrones, los electrones provenientes del NADH y FADH2 son transportados por citocromos. ¿Cuál molécula actúa como el "aceptor final" de esos electrones, convirtiéndose en agua?}
    \begin{enumerate}
        \item[A)] Oxígeno ($)
        \item[B)] Dióxido de carbono ($)
        \item[C)] Citocromo C
        \item[D)] ATP sintasa
    \end{enumerate}
    \textit{Respuesta: A. El oxígeno tiene gran electronegatividad y al recibir los electrones y protones, se reduce formando $. Sin $, la cadena se detiene (asfixia celular).}

    \item \textbf{En el metabolismo muscular humano bajo condiciones anaeróbicas (ejercicio intenso), el piruvato no puede entrar a la mitocondria. En cambio, es fermentado a:}
    \begin{enumerate}
        \item[A)] Etanol y $
        \item[B)] Acetil-CoA
        \item[C)] Ácido láctico (Lactato)
        \item[D)] Ácido cítrico
    \end{enumerate}
    \textit{Respuesta: C. La fermentación láctica oxida el NADH a NAD+ para que la glucólisis continúe produciendo 2 ATP. El etanol lo producen levaduras.}

    \item \textbf{¿Cuál es la función biológica de una enzima?}
    \begin{enumerate}
        \item[A)] Aumentar la energía de activación de una reacción.
        \item[B)] Disminuir la energía de activación, acelerando la velocidad de reacción.
        \item[C)] Alterar el equilibrio térmico para evitar el calentamiento celular.
        \item[D)] Consumirse permanentemente al catalizar una reacción.
    \end{enumerate}
    \textit{Respuesta: B. Las enzimas son catalizadores biológicos que aceleran las reacciones disminuyendo la barrera de energía requerida para iniciar, y no se consumen en el proceso.}

    \item \textbf{Un inhibidor farmacológico ocupa reversiblemente el sitio activo de una enzima impidiendo la unión del sustrato. ¿De qué tipo de inhibición enzimática hablamos?}
    \begin{enumerate}
        \item[A)] Irreversible
        \item[B)] No competitiva (alostérica)
        \item[C)] Competitiva
        \item[D)] Alostérica activadora
    \end{enumerate}
    \textit{Respuesta: C. Como "compite" por el mismo sitio activo estructural, es inhibición competitiva. La no competitiva actúa en otro sitio.}

    \item \textbf{Según el Dogma Central de la Biología Molecular, el proceso de copia del ADN a ARNm se denomina:}
    \begin{enumerate}
        \item[A)] Replicación
        \item[B)] Traducción
        \item[C)] Mutación
        \item[D)] Transcripción
    \end{enumerate}
    \textit{Respuesta: D. Transcribir es pasar a un lenguaje similar (nucleótidos a nucleótidos). Traducir es pasar a otro lenguaje (nucleótidos a aminoácidos).}

    \item \textbf{En una molécula de ADN bicatenario de doble hélice, si la cantidad de Adenina (A) es del 20\%, ¿cuál es el porcentaje de Citocina (C) según la ley de Chargaff?}
    \begin{enumerate}
        \item[A)] 20\%
        \item[B)] 30\%
        \item[C)] 40\%
        \item[D)] 80\%
    \end{enumerate}
    \textit{Respuesta: B. A se une a T (20\% + 20\% = 40\%). Queda 60\% para G+C, por ende C es el 30\% y G el 30\%.}

    \item \textbf{A diferencia del ADN, la molécula de ARN contiene la base nitrogenada:}
    \begin{enumerate}
        \item[A)] Timina
        \item[B)] Uracilo
        \item[C)] Guanina
        \item[D)] Adenina
    \end{enumerate}
    \textit{Respuesta: B. El ARN carece de Timina, la cual es reemplazada por Uracilo.}

    \item \textbf{¿En qué región de la célula eucariota ocurre la síntesis de proteínas (Traducción)?}
    \begin{enumerate}
        \item[A)] Núcleo
        \item[B)] Nucléolo
        \item[C)] Aparato de Golgi
        \item[D)] En los ribosomas del citoplasma y RER
    \end{enumerate}
    \textit{Respuesta: D. Los ribosomas son las "fábricas de proteínas", y traducen el ARNm ensamblando aminoácidos en el citosol o Retículo Endoplásmico Rugoso.}

    \item \textbf{El enlace principal que estabiliza la estructura primaria de una proteína es:}
    \begin{enumerate}
        \item[A)] El puente de hidrógeno
        \item[B)] El enlace disulfuro
        \item[C)] El enlace peptídico
        \item[D)] Interacciones hidrofóbicas
    \end{enumerate}
    \textit{Respuesta: C. La secuencia lineal de aminoácidos (1ra) se da por enlaces peptídicos (covalentes). Los puentes de H forman la secundaria (alfa hélice, hoja beta).}

    \item \textbf{Forma principal de almacenamiento de carbohidratos en el hígado y músculos de los vertebrados:}
    \begin{enumerate}
        \item[A)] Celulosa
        \item[B)] Almidón
        \item[C)] Glucógeno
        \item[D)] Quitina
    \end{enumerate}
    \textit{Respuesta: C. El glucógeno es una reserva de glucosa muy ramificada. El almidón y la celulosa son exclusivos de plantas. La quitina está en el exoesqueleto de artrópodos y pared de hongos.}

    \item \textbf{Molécula de alta energía (moneda energética universal celular) sintetizada mayormente en las mitocondrias:}
    \begin{enumerate}
        \item[A)] NAD+
        \item[B)] GTP
        \item[C)] ATP (Adenosín Trifosfato)
        \item[D)] FADH2
    \end{enumerate}
    \textit{Respuesta: C. El ATP almacena energía química en los enlaces fosfoanhídrido entre sus tres grupos fosfato.}

    \item \textbf{¿Qué enzima realiza la síntesis del ADN y posee actividad de corrección de errores durante la Fase S del ciclo celular?}
    \begin{enumerate}
        \item[A)] ARN polimerasa
        \item[B)] Helicasa
        \item[C)] Ligasa
        \item[D)] ADN polimerasa
    \end{enumerate}
    \textit{Respuesta: D. La ADN Polimerasa añade nucleótidos y revisa. La Helicasa rompe puentes de H, y la Ligasa une los fragmentos de Okazaki.}

    \item \textbf{Un codón consta de cuántos nucleótidos en el ARNm que codifican para un solo aminoácido?}
    \begin{enumerate}
        \item[A)] 1
        \item[B)] 2
        \item[C)] 3
        \item[D)] 4
    \end{enumerate}
    \textit{Respuesta: C. El código genético se lee en tripletes (3 nucleótidos = 1 codón = 1 aminoácido).}

    \item \textbf{El colesterol es el precursor bioquímico indispensable para la síntesis de:}
    \begin{enumerate}
        \item[A)] Proteínas de transporte sanguíneo
        \item[B)] Hormonas esteroideas (estrógenos, testosterona, cortisol)
        \item[C)] Insulina y glucagón
        \item[D)] Ácido Úrico
    \end{enumerate}
    \textit{Respuesta: B. El colesterol, siendo un lípido esteroideo, es la molécula base para formar hormonas sexuales y de la corteza suprarrenal. La insulina es peptídica.}

\end{enumerate}
\end{document}
